\newpage

\subsection{Messung der Thermischen Zeitkonstante des PT-1000}
\subsubsection{Aufgabenstellung}

\subsubsection{Verwendetes Material}

\begin{table}[h]
	\centering
	\caption{Verwendetes Material}
	\begin{tabularx}{\textwidth}{XXX}
		\textbf{Material} & \textbf{Verwendung} & \textbf{Anmerkungen} \\ \midrule
		Messverstärker & $A_1$, $A_2$, $R$, $R_p$ & Offsetspannung abgleichen! \\
		Temperatursensor & PT-1000 ($R + \Delta R(\vartheta)$) & $R=\SI{1}{\kilo\ohm}, \Delta R(\SI{0}{\degree}) = \SI{0}{\ohm}$ \\
		Labornetzteil & Hilfsspannung ($U_H$) & $U_H = \SI{-1}{\volt}$ \\
		& & \\
	\end{tabularx}	
\end{table}
\subsubsection{Versuchsaufbau}

Bevor die Schaltung aufgebaut wird, müssen jeweils die Offsetspannungen der beiden Operationsverstärker wie in \autoref{sec:1-1} beschrieben abgeglichen werden.

\begin{figure}[h]
	\centering
	\begin{circuitikz}[scale=0.9, transform shape]
		\draw (0,0) to[vR, l_=$R + \Delta R(\vartheta)$, o-o] (0,2) coordinate(pt1000)
			to[R, l_=$R$] ++(0,2) to[short, -*] ++(1.5,0) node[vcc]{$U_H = -1V$} -- ++(1.5,0)
			to[R=$R$, -*] ++(0,-2) coordinate(sense) to[R=$R$] ++(0,-2) -- ++(-3,0);
		\myNIA{A2}{(sense)}{$\SI{10}{\kilo\ohm} - R_P$}{$R_P$};
		\draw (0,0) node[op amp, anchor=out](A1){\texttt{A1}};
		\draw (A1.-) |- (pt1000);
		\draw (A1.+) -- (A1.+ |- A2-gnd) node[ground]{};
		\draw (A2-out) to[open, v=$U_A$, o-o] (A2-out |- A2-gnd) node[ground]{};
	\end{circuitikz}
	\caption{Messschaltung}
\end{figure}

\subsubsection{Berechnungen}

\subsubsection{Erkenntnisse}
